\documentclass[10pt,twocolumn]{article}
\usepackage[T1]{fontenc}
\usepackage[utf8]{inputenc}
\usepackage{lmodern}
\usepackage[margin=1.8cm,top=2.4cm,bottom=2cm,columnsep=0.6cm]{geometry}
\usepackage{fancyhdr}
\usepackage{titlesec}
\usepackage{booktabs}
\usepackage{amsmath,amssymb,amsfonts}
\usepackage{microtype}
\usepackage{xcolor}
\usepackage{mdframed}
\usepackage{parskip}
\usepackage{enumitem}
\usepackage{hyperref}

\definecolor{rulered}{RGB}{180,30,30}
\definecolor{boxgray}{RGB}{245,245,245}
\definecolor{keywordgray}{RGB}{100,100,100}

\pagestyle{fancy}
\fancyhf{}
\fancyhead[L]{\textsc{\small Claude Mag}}
\fancyhead[C]{\small\textit{Machine Learning Research Digest}}
\fancyhead[R]{\small 2026-10-12}
\fancyfoot[C]{\small\thepage}
\renewcommand{\headrulewidth}{0.8pt}

\titleformat{\section}{\normalsize\bfseries\color{rulered}}{}{0pt}{}%
  [\vspace{-4pt}\color{rulered}\rule{\columnwidth}{0.4pt}]
\titlespacing{\section}{0pt}{8pt}{4pt}

\hypersetup{colorlinks=true,urlcolor=rulered,linkcolor=black}

\begin{document}
\setlength{\parindent}{0pt}
\setlength{\parskip}{4pt}

{\small\color{keywordgray}\textbf{Keywords:}
  Bayesian neural networks \textbullet{}
  confidence sequences \textbullet{}
  adaptive inference \textbullet{}
  uncertainty quantification}\par\vspace{4pt}

{\LARGE\bfseries\raggedright
  When Has a Bayesian Neural Network Sampled Enough?
  Adaptive Inference Time with Statistical Guarantees}\par\vspace{4pt}

{\large\itshape\raggedright
  Anytime-Valid Confidence Sequences Terminate Monte Carlo Sampling
  Per-Input When the Decision Is Statistically Resolved}\par\vspace{6pt}

{\small\textbf{By} Fabian Denoodt \& Sibylle Hess}\par\vspace{2pt}

{\small\color{keywordgray}arXiv:2610.12212 \textbullet{} NeurIPS 2026 Workshop \textbullet{} October 2026}
\par\vspace{2pt}\noindent{\color{rulered}\rule{\columnwidth}{0.6pt}}\vspace{6pt}

A fixed Monte Carlo sample budget in Bayesian neural network inference
over-samples easy inputs and may under-sample hard ones, with no formal
control on decision reliability. This paper terminates sampling per-input
using \emph{confidence sequences}---anytime-valid probability bounds---as
soon as the relevant decision is statistically resolved, concentrating
compute on ambiguous cases and providing a formal per-prediction guarantee.

\begin{mdframed}[backgroundcolor=boxgray,linecolor=rulered,linewidth=1pt,
    innerleftmargin=6pt,innerrightmargin=6pt,
    innertopmargin=6pt,innerbottommargin=6pt]
\textbf{\small AT A GLANCE}\par\vspace{4pt}
\begin{description}[leftmargin=0pt,labelsep=4pt,itemsep=2pt,parsep=0pt,
    font=\small\bfseries]
  \item[Problem:] Fixed-budget BNN inference leaves finite-sample
    prediction error uncontrolled and wastes compute on easy inputs.
  \item[Why it matters:] Reliable per-input guarantees are critical
    for high-stakes BNN deployment.
  \item[Method:] Anytime-valid KL-mixture confidence sequences
    for argmax, distributional, and threshold BNN decisions.
  \item[Result:] Adaptive sampling reduces latency vs.\ matched-reliability
    fixed budgets while providing a formal guarantee per prediction.
  \item[Evidence:] Image classification benchmarks; comparison with
    fixed-$S$ and heuristic early stopping.
\end{description}
\end{mdframed}

\section{The Big Picture}

Bayesian neural networks approximate the posterior predictive distribution
by averaging over $S$ Monte Carlo weight samples:
\[
\hat{p}(y \mid \mathbf{x}) = \frac{1}{S}\sum_{s=1}^S p(y \mid \mathbf{x}, \mathbf{w}^{(s)})
\]
The choice of $S$ is arbitrary, uniform across inputs, and provides no
bound on the finite-sample error of the resulting prediction. For inputs
far from the decision boundary, $S=5$ is likely sufficient; for borderline
cases, $S=1000$ may still be insufficient. A method that terminates
sampling when the decision is resolved---with a provable guarantee---would
improve both reliability and efficiency simultaneously.

\section{Why Existing Approaches Fall Short}

\textbf{Fixed-$S$ inference} ignores input difficulty, over-sampling easy
inputs and providing no guarantee on hard ones.

\textbf{Heuristic early stopping} (stopping when the running mean
stabilises) provides no formal guarantee and empirically fails near
decision boundaries.

\textbf{Conformal prediction} provides marginal set-coverage guarantees
over the test set, not per-input anytime guarantees tied to the sampling
process.

\textbf{Sequential hypothesis testing} (SPRT) was designed for binary
tests; multi-class prediction and distributional estimation require
a different treatment.

\section{Core Mechanism}

A \textbf{confidence sequence} $\{C_s\}_{s=1}^\infty$ satisfies:
\[
\Pr\!\bigl[\exists s \geq 1 : \theta \notin C_s\bigr] \leq \alpha
\]
This anytime-valid property allows stopping at any data-dependent time
$\tau$ without inflating the Type~I error. Unlike classical confidence
intervals (valid only at a fixed $s$), confidence sequences shrink
as more samples arrive and remain valid throughout.

The paper instantiates confidence sequences for three BNN tasks:
(1)~argmax class identification, (2)~full predictive distribution
estimation, and (3)~probability-threshold decisions.

\section{Technical Formulation}

Let $\hat{p}_s^c = \frac{1}{s}\sum_{j=1}^s p(y=c \mid \mathbf{x}, \mathbf{w}^{(j)})$.
A KL-mixture confidence sequence gives an anytime interval $(L_s^c, U_s^c)$:
\[
\Pr\!\Bigl[\exists s \geq 1 : p^c \notin (L_s^c, U_s^c)\Bigr] \leq \alpha
\]

The stopping rule for argmax identification is:
\[
\tau^* = \min\!\Bigl\{s : U_s^{c^*} > \max_{c \neq c^*} U_s^c \Bigr\}
\]
where $c^* = \arg\max_c \hat{p}_s^c$.

For distributional estimation with tolerance $\varepsilon$:
\[
\tau^\varepsilon = \min\!\Bigl\{s : r_s \leq \varepsilon\Bigr\},
\quad r_s = O\!\left(\sqrt{\frac{\log s}{s}}\right)
\]

For threshold decisions at level $\tau_\text{thr}$:
\[
\tau^{\tau_\text{thr}} = \min\!\Bigl\{s : L_s^{c^*} > \tau_\text{thr}
\text{ or } U_s^{c^*} < \tau_\text{thr}\Bigr\}
\]

\section{Learning or Inference Procedure}

At inference time, for each input $\mathbf{x}$:
\begin{enumerate}[itemsep=2pt,parsep=0pt]
  \item Draw $\mathbf{w}^{(s)} \sim q(\mathbf{w} \mid \mathcal{D})$ and
    evaluate $p(y \mid \mathbf{x}, \mathbf{w}^{(s)})$.
  \item Update running mean $\hat{\mathbf{p}}_s$ and CS bounds
    $(L_s^c, U_s^c)$.
  \item Check the stopping criterion for the chosen task.
  \item If met, return the prediction with its guarantee; else draw $s+1$.
\end{enumerate}
No training modifications are required; the method wraps any BNN.

\section{What the Guarantee Says}

For any data-dependent stopping time $\tau$ induced by the confidence
sequence and any error level $\alpha$, the decision at time $\tau$
has error probability at most $\alpha$---regardless of the input
distribution or Monte Carlo sample realisations.

\section{Experimental Findings}

\begin{itemize}[itemsep=2pt,parsep=0pt]
  \item The method assigns \textbf{more samples to ambiguous inputs}
    (near decision boundaries) and fewer to easy inputs, naturally
    concentrating compute.
  \item \textbf{Overall inference latency is reduced} vs.\ a fixed-$S$
    budget that provides the same average reliability.
  \item CS bounds are \textbf{tight in practice}: average sample count
    closely tracks input difficulty measured by entropy.
  \item For threshold decisions, unnecessary high-sample expenditure
    on clearly low- or high-uncertainty inputs is avoided.
\end{itemize}

\section{Ablations and Interpretation}

\textbf{Confidence level $\alpha$:} Stricter guarantees (lower $\alpha$)
require more samples per input; the relationship is approximately
$O(\log(1/\alpha))$.

\textbf{Posterior quality:} Better-calibrated posteriors tighten CS bounds
faster; an over-dispersed posterior still provides a valid guarantee
but at greater sample cost.

\textbf{Heuristic early stopping:} Provides no guarantee and fails
empirically on high-entropy inputs where the running mean converges slowly.

\textbf{Dataset difficulty skew:} Savings are largest on datasets with
many easy and few hard examples---the common case in practice.

\section*{Reference}

Fabian Denoodt, Sibylle Hess.
\textbf{When Has a Bayesian Neural Network Sampled Enough?\ Adaptive
Inference Time with Statistical Guarantees}.
NeurIPS 2026 Workshop (E-Values: From Statistics to ML).
arXiv:2610.12212, October 2026.\\
\url{https://arxiv.org/abs/2610.12212}

\end{document}
